% Generated draft: requires mathematical and source review.
\documentclass[11pt]{article}
\usepackage[margin=25mm]{geometry}
\usepackage{fontspec}
\setmainfont{Latin Modern Roman}
\usepackage{amsmath,amssymb,mathtools,unicode-math}
\setmathfont{Latin Modern Math}
\usepackage{enumitem,needspace,xurl,hyperref}
\setlength{\parindent}{0pt}
\setlength{\parskip}{5pt}
\setlength{\emergencystretch}{3em}
\newcommand{\N}{\mathbb{N}}
\newcommand{\Z}{\mathbb{Z}}
\newcommand{\Q}{\mathbb{Q}}
\newcommand{\R}{\mathbb{R}}
\newcommand{\C}{\mathbb{C}}
\newcommand{\F}{\mathbb{F}}
\newcommand{\A}{\mathbb{A}}
\newcommand{\PP}{\mathbb{P}}
\newcommand{\E}{\mathbb{E}}
\newcommand{\eps}{\varepsilon}
\newcommand{\dd}{\,\mathrm{d}}
\begin{document}
\section*{UnsolvedMath ID: 30005127}
\textbf{Source problem code:} catalogue entry
\subsection{Definitions and mathematical statement}
The supplied record for permanent UnsolvedMath ID \(30005127\), source code \texttt{catalogue entry}, has the title ``Open Problems in K-Theory of Ring Pushouts.'' Its statement is marked \texttt{unrecoverable}, and no clean mathematical statement is supplied.

The recorded passage is a table-of-contents fragment containing, in order, the word ``conjecture,'' dot leaders and page number \(1354\), followed by ``Georg Tamme (joint with Markus Land),'' the title ``K-theory of pushouts of associative rings,'' and dot leaders and page number \(1357\). The record adds that this passage is retained verbatim as context for the extracted problem.

\textbf{Unresolved mathematical statement.} No conjectural assertion, definitions, hypotheses, domains, boundary conditions, or quantifiers can be recovered from this fragment. In particular, the fragment does not establish that the preceding word ``conjecture'' belongs to the contribution about pushouts of associative rings, nor does it identify a conjecture attributable to Georg Tamme or Markus Land. The supplied problem title alone does not determine a mathematical claim. Consequently, no precise mathematical proposition can be stated from the supplied evidence without inventing content.

The registry labels this record ``open.'' Its supplied assessment explains that this classification follows an uncertainty rule because no definite claim can be identified; it is not verification that a particular mathematical conjecture remains open.
\subsection{Short English statement}
This record preserves a table-of-contents fragment, not an identifiable conjecture. It names a contribution by Georg Tamme, jointly with Markus Land, about K-theory of pushouts of associative rings, but provides no mathematical assertion to formalize. Recovering the intended problem requires further source evidence and human mathematical review.
\subsection{Sources}
\begin{itemize}
\item Supplied local UnsolvedMath registry record, permanent ID \(30005127\), source code \texttt{catalogue entry}, in the ``Oberwolfach Reports Open Problems'' collection. The record supplies the title ``Open Problems in K-Theory of Ring Pushouts,'' status ``open,'' statement status \texttt{unrecoverable}, and no clean statement.
\item Supplied report citation: \emph{Algebraic K-Theory (2023)}, Oberwolfach Reports, DOI \texttt{10.4171/owr/2022/24}, pp.\ 1313--1376. Supplied source and literature-source link, labeled ``Oberwolfach report'': \url{https://doi.org/10.4171/owr/2022/24}. This citation identifies the reported provenance, but the supplied fragment contains no recoverable authoritative mathematical assertion.
\item Attribution in the fragment: Georg Tamme, ``K-theory of pushouts of associative rings,'' joint with Markus Land, with table-of-contents page number \(1357\). The preceding ``conjecture'' fragment has page number \(1354\). No proposer of the extracted problem is supplied.
\item Supplied literature assessment dated exactly \texttt{2026-08-22}: no definite mathematical claim can be identified from the table-of-contents fragment, so the record is classified as open under the supplied uncertainty rule. This is supplied dated curation, not an independently verified literature assessment or source visit. The supplied background requests independent verification of time-sensitive status and cited claims.
\item The record supplies proposed year \(2022\) and the tag \texttt{owr-year:2022}; the report citation uses \(2023\). These year labels are retained without resolving their relationship.
\end{itemize}
\end{document}
